%% ============================================================
\section{The Weyl Lagrangian}
%% ============================================================

\subsection{Lagrangian terms}

For a single left-handed Weyl field $\psi_a(x)$, the Lorentz-invariant,
hermitian Lagrangian quadratic in $\psi$ is (eq.~36.2):
\begin{equation}
  \boxed{
    \mathcal{L}
    = i\psi^\dagger_{\dot{a}}\,\bar{\sigma}^{\mu\,\dot{a}\alpha}\,\partial_\mu\psi_\alpha
    - \tfrac{1}{2}m\,\varepsilon^{\alpha\beta}\psi_\alpha\psi_\beta
    - \tfrac{1}{2}m\,\varepsilon_{\dot{\alpha}\dot{\beta}}\psi^{\dagger\dot{\alpha}}\psi^{\dagger\dot{\beta}}
  }
  \tag{36.2}
\end{equation}

\medskip
\noindent\textbf{Cadabra2 expressions:}
\begin{align*}
  \mathcal{L}_{\text{kin}} &= i \psi^{\dagger \dal} \sigmabar^{\mu}\,_{\dal \alpha} \partial_{\mu}\left(\psi^{\alpha}\right) \\
  \mathcal{L}_{\text{mass}} &=  - \frac{1}{2}\,m \epsilon^{\alpha \beta} \psi_{\alpha} \psi_{\beta} \\
  \mathcal{L}_{\text{mass}}^\dagger &=  - \frac{1}{2}\,m \epsilon_{\dal \dbe} \psi^{\dagger \dal} \psi^{\dagger \dbe}
\end{align*}

\subsection{Hermiticity of the kinetic term}

Although $i\psi^\dagger\bar{\sigma}^\mu\partial_\mu\psi$ is not individually
hermitian, it satisfies (eq.~36.1):
\begin{equation}
  (i\psi^\dagger\bar{\sigma}^\mu\partial_\mu\psi)^\dagger
  = i\psi^\dagger\bar{\sigma}^\mu\partial_\mu\psi
    - i\partial_\mu(\psi^\dagger\bar{\sigma}^\mu\psi)\,.
  \tag{36.1}
\end{equation}
The second term is a total divergence and vanishes in the action.
Thus the kinetic term \emph{is} real in $S = \int d^4x\,\mathcal{L}$.

\subsection{Mass term and Majorana mass}

The mass term $\varepsilon^{\alpha\beta}\psi_\alpha\psi_\beta = \psi^\alpha\psi_\alpha$
is a Lorentz scalar; it is nonzero because $\psi_\alpha$ are Grassmann-valued
(fermionic) fields.  The parameter $m$ is a ``Majorana mass'':
\begin{itemize}
  \item The phase of $m$ is unphysical: redefining $\psi \to e^{-i\alpha/2}\psi$
        removes it.  We therefore take $m$ real and positive.
  \item There is no conserved U(1) charge associated with this mass term---it
        explicitly breaks any $\psi\to e^{i\alpha}\psi$ symmetry.
\end{itemize}

%% ============================================================
\section{Equations of Motion}
%% ============================================================

Varying the action $S = \int d^4x\,\mathcal{L}$ with respect to
$\psi^{\dagger\dot{a}}$ (holding $\psi$ fixed):
\begin{equation}
  0 = -\frac{\delta S}{\delta\psi^{\dagger\dot{a}}}
    = -i\,\bar{\sigma}^{\mu\,\dot{a}c}\,\partial_\mu\psi_c
      + m\,\psi^{\dagger\dot{a}}\,.
  \tag{36.3--36.4}
\end{equation}

\noindent In Cadabra2: $-\,i \sigmabar^{\mu}\,_{\dal \alpha} \partial_{\mu}\left(\psi^{\alpha}\right) + m \psi^{\dagger \dal}$

\medskip
\noindent The hermitian conjugate equation (varying w.r.t.\ $\psi^\alpha$):
\begin{equation}
  0 = -i\,\sigma^\mu_{a\dot{c}}\,\partial_\mu\psi^{\dagger\dot{c}}
      + m\,\psi_a\,.
  \tag{36.5}
\end{equation}

\noindent In Cadabra2: $-\,i \sigma^{\mu}\,_{\alpha \dal} \partial_{\mu}\left(\psi^{\dagger \dal}\right) + m \psi_{\alpha}$

\subsection{Combined 4-component form}

Equations~(36.4) and~(36.5) combine into the matrix equation (eq.~36.6):
\begin{equation}
  \begin{pmatrix}
    m\,\delta_a{}^c & -i\,\sigma^\mu_{a\dot{c}}\,\partial_\mu \\
    -i\,\bar{\sigma}^{\mu\,\dot{a}c}\,\partial_\mu & m\,\delta^{\dot{a}}{}_{\dot{c}}
  \end{pmatrix}
  \begin{pmatrix} \psi_c \\ \psi^{\dagger\dot{c}} \end{pmatrix}
  = 0\,.
  \tag{36.6}
\end{equation}
The $2\times 2$ diagonal blocks are proportional to $m\cdot\mathbf{1}_2$;
the off-diagonal blocks mix the left- and right-handed sectors via $\sigma^\mu$
and $\bar{\sigma}^\mu$.

%% ============================================================
\section{Gamma Matrices in the Weyl Representation}
%% ============================================================

\subsection{Definition}

The gamma matrices are defined by the block structure (eq.~36.7):
\begin{equation}
  \gamma^\mu \equiv
  \begin{pmatrix}
    0 & \sigma^\mu_{a\dot{c}} \\
    \bar{\sigma}^{\mu\,\dot{a}c} & 0
  \end{pmatrix}.
  \tag{36.7}
\end{equation}
This is the \textbf{Weyl (chiral) representation}.

\subsection{Explicit matrices}

\begin{align}
  \gamma^0 &= \begin{pmatrix}0 & 0 & 1 & 0 \\ 0 & 0 & 0 & 1 \\ 1 & 0 & 0 & 0 \\ 0 & 1 & 0 & 0\end{pmatrix},&
  \gamma^1 &= \begin{pmatrix}0 & 0 & 0 & 1 \\ 0 & 0 & 1 & 0 \\ 0 & -1 & 0 & 0 \\ -1 & 0 & 0 & 0\end{pmatrix}\\[6pt]
  \gamma^2 &= \begin{pmatrix}0 & 0 & 0 & -i \\ 0 & 0 & i & 0 \\ 0 & i & 0 & 0 \\ -i & 0 & 0 & 0\end{pmatrix},&
  \gamma^3 &= \begin{pmatrix}0 & 0 & 1 & 0 \\ 0 & 0 & 0 & -1 \\ -1 & 0 & 0 & 0 \\ 0 & 1 & 0 & 0\end{pmatrix}
\end{align}

\noindent Written symbolically in terms of $\sigma^\mu = (I, \vec{\sigma})$
and $\bar{\sigma}^\mu = (I,-\vec{\sigma})$:

\[
  \gamma^0 = \begin{pmatrix}0&I\\I&0\end{pmatrix},\quad
  \gamma^k = \begin{pmatrix}0&\sigma_k\\-\sigma_k&0\end{pmatrix}
  \quad (k=1,2,3)\,.
\]

%% ============================================================
\section{Clifford Algebra Verification}
%% ============================================================

\subsection{The algebra}

The gamma matrices obey the Clifford algebra (eq.~36.9):
\begin{equation}
  \boxed{\{\gamma^\mu,\,\gamma^\nu\} \equiv \gamma^\mu\gamma^\nu + \gamma^\nu\gamma^\mu
  = -2\,g^{\mu\nu}\,I_4\,,}
  \tag{36.9}
\end{equation}
with $g^{\mu\nu}=\mathrm{diag}(-1,+1,+1,+1)$ (Srednicki mostly-plus convention).

This follows from the sigma-matrix identities (eq.~36.8):
\begin{align}
  \sigma^\mu\bar{\sigma}^\nu + \sigma^\nu\bar{\sigma}^\mu &= -2g^{\mu\nu}\,I_2\,,
  \tag{36.8a}\\
  \bar{\sigma}^\mu\sigma^\nu + \bar{\sigma}^\nu\sigma^\mu &= -2g^{\mu\nu}\,I_2\,.
  \tag{36.8b}
\end{align}
\verified{$\sigma^\mu\bar{\sigma}^\nu+\sigma^\nu\bar{\sigma}^\mu=-2g^{\mu\nu}I_2$: max error $0.0e+00$.}

\subsection{Numerical verification table}

All ten independent anticommutator pairs are verified below:
\medskip
\begin{center}
\begin{tabular}{@{}llll@{}}
  \toprule
  Pair & Formula & Value & Verified \\
  \midrule
  $\{\gamma^0,\gamma^0\}$ & $-2g^{00}$ & $+2 \cdot I_4$ & \textcolor{verifygreen}{\checkmark} \\
  $\{\gamma^0,\gamma^1\}$ & $-2g^{01}$ & $-0 \cdot I_4$ & \textcolor{verifygreen}{\checkmark} \\
  $\{\gamma^0,\gamma^2\}$ & $-2g^{02}$ & $-0 \cdot I_4$ & \textcolor{verifygreen}{\checkmark} \\
  $\{\gamma^0,\gamma^3\}$ & $-2g^{03}$ & $-0 \cdot I_4$ & \textcolor{verifygreen}{\checkmark} \\
  $\{\gamma^1,\gamma^1\}$ & $-2g^{11}$ & $-2 \cdot I_4$ & \textcolor{verifygreen}{\checkmark} \\
  $\{\gamma^1,\gamma^2\}$ & $-2g^{12}$ & $-0 \cdot I_4$ & \textcolor{verifygreen}{\checkmark} \\
  $\{\gamma^1,\gamma^3\}$ & $-2g^{13}$ & $-0 \cdot I_4$ & \textcolor{verifygreen}{\checkmark} \\
  $\{\gamma^2,\gamma^2\}$ & $-2g^{22}$ & $-2 \cdot I_4$ & \textcolor{verifygreen}{\checkmark} \\
  $\{\gamma^2,\gamma^3\}$ & $-2g^{23}$ & $-0 \cdot I_4$ & \textcolor{verifygreen}{\checkmark} \\
  $\{\gamma^3,\gamma^3\}$ & $-2g^{33}$ & $-2 \cdot I_4$ & \textcolor{verifygreen}{\checkmark} \\
  \bottomrule
\end{tabular}
\end{center}
\medskip
\verified{Clifford algebra: all 10 pairs verified. Max error: $0.0e+00$.}

\subsection{The \texorpdfstring{$\gamma^5$}{gamma5} matrix}

Define (eq.~36.46):
\begin{equation}
  \gamma^5 \equiv i\gamma^0\gamma^1\gamma^2\gamma^3
  = \begin{pmatrix}-I_2 & 0 \\ 0 & +I_2\end{pmatrix}.
  \tag{36.43}
\end{equation}
Numerically:
\[
  \gamma^5 = \begin{pmatrix}-1 & 0 & 0 & 0 \\ 0 & -1 & 0 & 0 \\ 0 & 0 & 1 & 0 \\ 0 & 0 & 0 & 1\end{pmatrix}\,.
\]
\verified{$\gamma^5 = \mathrm{diag}(-I_2,+I_2)$: error $0.0e+00$.}

\medskip
\noindent The left- and right-handed \textbf{projection operators} are (eq.~36.44):
\begin{align}
  P_L &\equiv \tfrac{1}{2}(I - \gamma^5)
     = \begin{pmatrix}I_2&0\\0&0\end{pmatrix},
  &
  P_R &\equiv \tfrac{1}{2}(I + \gamma^5)
     = \begin{pmatrix}0&0\\0&I_2\end{pmatrix}.
  \tag{36.44}
\end{align}
\verified{$P_L^2=P_L$, $P_R^2=P_R$, $P_LP_R=0$: max error $0.0e+00$.}

%% ============================================================
\section{Majorana 4-Component Spinor}
%% ============================================================

\subsection{Definition}

The \textbf{Majorana spinor} is built from a single left-handed Weyl field $\psi_a$
(eq.~36.10):
\begin{equation}
  \Psi \equiv \begin{pmatrix}\psi_c\\\psi^{\dagger\dot{c}}\end{pmatrix}.
  \tag{36.10}
\end{equation}
Equation~(36.6) then becomes the \textbf{Dirac equation} (eq.~36.11):
\begin{equation}
  \boxed{(-i\gamma^\mu\partial_\mu + m)\Psi = 0\,.}
  \tag{36.11}
\end{equation}

\subsection{Majorana condition}

The Majorana spinor is its own \textbf{charge conjugate}:
\begin{equation}
  \Psi^C \equiv \mathcal{C}\,\overline{\Psi}^{\,T} = \Psi\,,
\end{equation}
where the charge conjugation matrix is (eq.~36.32):
\begin{equation}
  \mathcal{C} = \begin{pmatrix}\varepsilon_{ac}&0\\0&\varepsilon^{\dot{a}\dot{c}}\end{pmatrix}
  = \begin{pmatrix}0 & -1 & 0 & 0 \\ 1 & 0 & 0 & 0 \\ 0 & 0 & 0 & 1 \\ 0 & 0 & -1 & 0\end{pmatrix}\,.
  \tag{36.32}
\end{equation}

The matrix $\mathcal{C}$ satisfies (eq.~36.36):
\begin{equation}
  \mathcal{C}^T = \mathcal{C}^\dagger = \mathcal{C}^{-1} = -\mathcal{C}\,.
  \tag{36.36}
\end{equation}
and conjugates the gamma matrices as (eq.~36.40):
\begin{equation}
  \mathcal{C}^{-1}\gamma^\mu\mathcal{C} = -(\gamma^\mu)^T\,.
  \tag{36.40}
\end{equation}
\verified{$\mathcal{C}^{-1}\gamma^\mu\mathcal{C}=-(\gamma^\mu)^T$
for all $\mu$: max error $0.0e+00$.}

\subsection{Majorana Lagrangian}

Re-expressing eq.~(36.2) using $\overline{\Psi}=\Psi^\dagger\beta$ (where $\beta=\gamma^0$
numerically) gives (eq.~36.41):
\begin{equation}
  \mathcal{L} = \tfrac{i}{2}\overline{\Psi}\gamma^\mu\partial_\mu\Psi
              - \tfrac{1}{2}m\,\overline{\Psi}\Psi\,.
  \tag{36.41}
\end{equation}

%% ============================================================
\section{Dirac Fermion from Two Weyl Fields}
%% ============================================================

\subsection{Two-Weyl-field Lagrangian}

Begin with two left-handed fields $\chi_a$, $\xi_a$ (rewritten from $\psi_{1,2}$
via eq.~36.14--36.15). The Lagrangian is (eq.~36.16):
\begin{equation}
  \mathcal{L}
  = i\chi^\dagger\bar{\sigma}^\mu\partial_\mu\chi
  + i\xi^\dagger\bar{\sigma}^\mu\partial_\mu\xi
  - m\,\chi\xi - m\,\xi^\dagger\chi^\dagger\,,
  \tag{36.16}
\end{equation}
invariant under the U(1) symmetry (eq.~36.17):
\begin{equation}
  \chi \to e^{-i\alpha}\chi\,, \qquad \xi \to e^{+i\alpha}\xi\,.
  \tag{36.17}
\end{equation}

\subsection{Dirac spinor and Lagrangian}

Define the \textbf{Dirac spinor} (eq.~36.19):
\begin{equation}
  \Psi \equiv \begin{pmatrix}\chi_c\\\xi^{\dagger\dot{c}}\end{pmatrix}\,,
  \qquad
  \overline{\Psi} = \Psi^\dagger\beta = (\xi^a,\,\chi^\dagger_{\dot{a}})\,.
  \tag{36.19, 36.22}
\end{equation}
Then the Dirac Lagrangian (eq.~36.28):
\begin{equation}
  \boxed{\mathcal{L} = i\overline{\Psi}\gamma^\mu\partial_\mu\Psi - m\overline{\Psi}\Psi\,,}
  \tag{36.28}
\end{equation}
where (eq.~36.23):
\begin{equation}
  \overline{\Psi}\Psi = \xi^a\chi_a + \chi^\dagger_{\dot{a}}\xi^{\dagger\dot{a}}\,.
  \tag{36.23}
\end{equation}

The Noether current for the U(1) symmetry (eq.~36.30):
\begin{equation}
  j^\mu = \overline{\Psi}\gamma^\mu\Psi
        = \chi^\dagger\bar{\sigma}^\mu\chi - \xi^\dagger\bar{\sigma}^\mu\xi\,.
  \tag{36.30}
\end{equation}

%% ============================================================
\section{Dirac vs.\ Majorana --- Comparison}
%% ============================================================

\begin{center}
\renewcommand{\arraystretch}{1.5}
\begin{tabular}{@{}lll@{}}
  \toprule
  \textbf{Property} & \textbf{Majorana} & \textbf{Dirac} \\
  \midrule
  Weyl fields & 1 (single $\psi$) & 2 ($\chi$ and $\xi$) \\
  U(1) symmetry & No & Yes ($\chi\to e^{-i\alpha}\chi$, $\xi\to e^{+i\alpha}\xi$) \\
  Charge & None (Majorana mass) & Conserved $Q=\int j^0 d^3x$ \\
  $\Psi^C = \Psi$ & Yes & No \\
  Lagrangian & $\frac{i}{2}\bar{\Psi}\gamma^\mu\partial_\mu\Psi - \frac{m}{2}\bar{\Psi}\Psi$ &
               $i\bar{\Psi}\gamma^\mu\partial_\mu\Psi - m\bar{\Psi}\Psi$ \\
  Factor of $\frac{1}{2}$ & Yes (overcounting) & No \\
  Scalar analogue & Real scalar $\phi=\phi^\dagger$ & Complex scalar $\phi\neq\phi^\dagger$ \\
  Neutrino model & Seesaw mechanism & Dirac neutrino \\
  Degrees of freedom & 2 (real) & 4 (complex) \\
  \bottomrule
\end{tabular}
\end{center}

%% ============================================================
\section{Summary}
%% ============================================================

\begin{center}
\renewcommand{\arraystretch}{1.4}
\begin{tabular}{@{}ll@{}}
  \toprule
  Object & Expression \\
  \midrule
  Weyl Lagrangian & $\mathcal{L}=i\psi^\dagger\bar\sigma^\mu\partial_\mu\psi
    -\frac{1}{2}m\psi\psi - \frac{1}{2}m\psi^\dagger\psi^\dagger$ \\
  EOM ($\delta/\delta\psi^\dagger$) & $-i\bar\sigma^{\mu\dot{a}c}\partial_\mu\psi_c
    +m\psi^{\dagger\dot{a}}=0$ \\
  $\gamma$ matrix (Weyl rep) & $\gamma^\mu=\bigl[\begin{smallmatrix}0&\sigma^\mu\\
    \bar\sigma^\mu&0\end{smallmatrix}\bigr]$ \\
  Clifford algebra & $\{\gamma^\mu,\gamma^\nu\}=-2g^{\mu\nu}I_4$ \\
  $\gamma^5$ & $i\gamma^0\gamma^1\gamma^2\gamma^3
    =\mathrm{diag}(-I_2,+I_2)$ \\
  Majorana spinor & $\Psi=(\psi_c,\psi^{\dagger\dot{c}})^T$,\quad $\Psi^C=\Psi$ \\
  Dirac equation & $(-i\gamma^\mu\partial_\mu+m)\Psi=0$ \\
  Dirac spinor & $\Psi=(\chi_c,\xi^{\dagger\dot{c}})^T$,\quad $\Psi^C\neq\Psi$ \\
  Dirac Lagrangian & $i\bar\Psi\gamma^\mu\partial_\mu\Psi - m\bar\Psi\Psi$ \\
  Charge conjugation & $\mathcal{C}^{-1}\gamma^\mu\mathcal{C}=-(\gamma^\mu)^T$ \\
  \bottomrule
\end{tabular}
\end{center}

\bigskip
\noindent\textbf{Next:} Chapter 37 develops the quantization of
Dirac and Majorana fields, introducing anticommutation relations for
the creation/annihilation operators.
